% !PRINTED 447
\weberpage{447}
%<<CONT
eight-fold. Since it remains unchanged by the substitutions of degree 2 $\omega$, $\Theta\omega$, $\Theta^2\omega$, $\Theta^3\omega$, and since the poles of these substitutions lie on $A$ (because $A$ remains unchanged by them), while $a$ does not lie on $A$, the principal axes of these four substitutions $\omega$, $\Theta\omega$, $\Theta^2\omega$, $\Theta^3\omega$ must pass through the point $a$.

The poles of the four substitutions $\omega$, $\Theta\omega$, $\Theta^2\omega$, $\Theta^3\omega$, which we shall call $a_1$, $a_2$, $a_3$, $a_4$, are, as images of the point $a$, likewise eight-fold poles, and must, like $a$, be poles of substitutions of the fourth order, all of which are in any case different from $\Theta$ and $\Theta^3$, because otherwise $\Theta^2$ would have to be equal to one of those four substitutions, which is not the case. Under $\omega$ only two points of the axis $A$, namely the poles of $\omega$ and of $\Theta^2\omega$, remain unchanged.

Let us now consider the two other poles $c_1$, $c_2$ of the substitution $\Theta$. These points likewise lie on the principal axis $A$, but are different from the points $a_1$, $a_2$, $a_3$, $a_4$, because these, as already remarked, do not occur among the poles of $\Theta$. Now $\omega\Theta\omega = \Theta$, hence

\[
\omega\,\Theta(x) = \Theta\,\omega(x),
\]
and if now $(x) = \Theta(x)$, then also

\[
\omega(x) = \Theta\,\omega(x);
\]
that is, if $(x)$ are the coordinates of a pole of $\Theta$, then the $\omega(x)$ have the same meaning. Under the transformation $\omega$, therefore, the poles of $\Theta$ merely pass into one another. The pole $a$ remains unchanged under $\omega$; $c_1$ and $c_2$, on the other hand, cannot remain unchanged, because they do not occur among the poles of $\omega$, and must therefore pass into one another.

We thus have the following overview:

\begin{center}
\begin{tabular}{@{}lllllll@{}}
Unchanged under & $\omega$ and $\Theta^2\omega$ & on & $A$ & only the points & $a_1, a_3$ & \\
\textquotedbl\ \textquotedbl & $\Theta\omega$\ \textquotedbl\ $\Theta^3\omega$ & \textquotedbl & $A$ & \textquotedbl\ \textquotedbl\ \textquotedbl & $a_2, a_4$ & \\
\textquotedbl\ \textquotedbl & $\Theta$\ \textquotedbl\ $\Theta^3$ & \textquotedbl & $A$ & \textquotedbl\ \textquotedbl\ \textquotedbl & $c_1, c_2.$ &
\end{tabular}
\end{center}

A point $x$ of $A$ different from these six passes into itself only under the substitutions 1 and $\Theta^2$, and may therefore be called a \gesperrt{two-fold pole}. Then the following theorem holds:

\noindent 3.\ \gesperrt{A two-fold pole passes, under the group of degree 8 derived from $\omega$ and $\Theta$, into four different positions.}
